\documentclass[10pt]{article}

\usepackage[utf8]{inputenc}

\usepackage[T1]{fontenc}

\usepackage{amsmath}

\usepackage{amsfonts}

\usepackage{amssymb}

\usepackage[version=4]{mhchem}

\usepackage{stmaryrd}

\usepackage{graphicx}

\usepackage[export]{adjustbox}

\graphicspath{ {./images/} }

\usepackage{bbold}



\begin{document}

Р. многочленов $f(x)$ и $g(x)$ с числовыми коәффиииентами можно представить в виде определителя порядка $n$ (или $s$ ). Для этого находят остаток от деления $x^{k} g(x)$ на $f(x), k=0,1,2, \ldots, n-1$. Пусть это будет



Тогда



\[

a_{k 0}+a_{k 1} x+\ldots+a_{k n-1} x^{n-1}

\]



\begin{center}

\includegraphics[max width=\textwidth]{2023_07_08_bdc971b48298cca3e951g-1}

\end{center}



Дискриминант $D(f)$ многочлена



\[

f(x)=a_{0} x^{n}+a_{1} x^{n-1} \ldots a_{n}, a_{0} \neq 0,

\]



выражается через Р. многочлена $f(x)$ ш его производной $f^{\prime}(x)$ следующим образом:



\[

D(f)=(-1)^{\frac{n(n-1)}{2}} a_{0}^{-1} R\left(f, f^{\prime}\right) .

\]



П ри и енение к ре шению у р а в н е н и й. Пусть дана система двух алгебраич. уравнений с коэффициентами из поля $P$ :



\[

\left.\begin{array}{l}

f(x, y)=0 \\

g(x, y)=0

\end{array}\right\}

\]



Многочлены $f$ и $g$ записывают по степеням $x$ :



\[

\begin{aligned}

& f(x, y)=a_{0}(y) x^{k}+a_{1}(y) x^{k-1}+\ldots+a_{k}(y), \\

& g(x, y)=b_{0}(y) x^{l}+b_{1}(y) x^{l-1}+\ldots+b_{l}(y),

\end{aligned}

\]



Эффективность замены кручения Уайтхеда на Р. к. основывается на т е о р е м е Б а с с а [4]: если $\pi-$ конечная группа, то элемент $\omega \in W h(\pi)$ имеет конечный порядок, если $h_{*}(\omega)=1$ для любого представления $h$, где $h_{*}(\omega)-\mathrm{P}$. к., индуци рованное элементом $\omega$.



Jum.: [1] R e i d e m e i s t e r K., "Abhandl. math. Semin.



\includegraphics[max width=\textwidth, center]{2023_07_08_bdc971b48298cca3e951g-1(2)}

d e «Матем. сб.», 1936, т. 1, № 5, c. 737-43; [4] B a s s H. "Publ. math.", 1964, №'22, p.' 5-60 [IHES). A.C.Мищенко.



РЕИНОЛЬДСА ЧИСЛО - один из критериев подобия для течений вязких жидкостей и газов, характеризующий соотношение между инерционными силами и силами вязкости:



\[

\operatorname{Re}=\rho v l / \mu,

\]



где $\rho$ - плотность, $\mu$ - динамич. коэффициент вязкости жидкости или газа, $v$ - характерная скорость потока, $l$ - характерный линейный размер.



От Р. ч. зависит также режим течения жидкости, характеризуемый к р и т и ч е с к и м Р.ч. Re кр. При $\operatorname{Re}<\operatorname{Re}_{\text {кр }}$ возможно лишь ламинарное течение жидкости, a при $\operatorname{Re}>\operatorname{Re}_{\text {кр }}$ течение может стать турбулентным.



Р. ч. названо по имени О. Рейнольдса (O. Reynolds).



По материалам одноименной статъи из БСЭ-з.



РЕКУ РРЕНТНАЯ ТОЧКА Д И н а м И ч е с К о Й с и с т е м - точка $x$ динамич. системы $f^{t}$ (или, в иных обозначениях $f(t, \cdot)$, см. [2]), заданной на метрич. пространстве $S$, удовлетворяющая условию: для всякого $\varepsilon>0$ найдется $T>0$ такое, что все точки траектории $f^{t} x$ содержатся в $\varepsilon$-окрестности всякой дуги временно́й длины $T$ этой траектории (иными словами,



\begin{center}

\includegraphics[max width=\textwidth]{2023_07_08_bdc971b48298cca3e951g-1(1)}

\end{center}



\[

\left\{f^{t} x\right\}, t \in[\tau, \tau+T]

\]



Говорят, что многочлен $F(y)$ получен путем исключения $x$ из многочленов $f(x, y)$ и $g(x, y)$. Если $x=\alpha, y=\beta-$ решение системы $(5)$, то $F(\beta)=0$, и обратно, если $F(\beta)=0$, то или многочлены $f(x, \beta), g(x, \beta)$ имеют общий корень (к-рый надо искать как корень их наибольшего общего делителя), или $a_{0}(\beta)=b_{0}(\beta)=0$. Тем самым решение системы (5) сводится к вычислению корней многочлена $F(y)$ и общих корней многочленов $f(x, \beta), g(x, \beta)$ с одним неизвестным.



Аналогично можно решать и системы уравнений с любым числом неизвестных, но эта задача приводит к весьма громоздким вычисл ниям (см. также Исключения теория).



Лum.: [1] К у р о ш А. Г., Курс высшей алгебры, 11 изд., Ј., 1949; [3] В а н д е р в а р п е н Б. J., Алггера, пер. с нем., 2 изд., М., 1979; [4] Х о д ж В., П и д о Д., Методы алІ ебраической геометрии, пер. с англ., т. 1-3, М., 1954-55.



РЕЙДЕМЕЙСТЕРА КРУЧЕНИЕ, К. К Р уч е н И д е Р а м а, к р у ч е н и е Ф р а ц ц а,- инвариант, позволяющий различать многие структуры в дифференциальной топологии, напр. узлы, гладкие структуры на многообразиях, в частности на линзовых пространствах. Впервые Р. к. введено К. Рейдемейстером (см. [1]) при изучении трехмерных линз, обобщения для $n$-мерных линз были независимо получены в [2] и [3].



Пусть $C$ - свободный комплекс левых $A$-модулей, где $A$ - ассоциативное кольцо с единицей. Пусть, далее, $h$ - матричное представление кольца $A$, т. е. гомоморфизм кольца $A$ в кольцо $R n \times n$ всех действительных $(n \times n)$-матриц. И пусть в модулях $C_{k}$ комплекса $C$ отмечены базисы $c_{k}$, а комплекс $C^{\prime}=R^{n \times n} \otimes_{A} C R^{n \times n_{-}}$ модулей ацикличен; тогда определено Уайтxеда кручение $\tau\left(C^{\prime}\right) \in \bar{K}_{1} R^{n \times n}=\bar{K}_{1} R=R_{+}$, где $R_{+}-$мультипликативная группа поля действительных чисел. Число $\tau\left(C^{\prime}\right)$ наз. к р у че н и е м Р е й де ме й с те р а комплекса $C^{\prime}$, а также дей с т в и тельны м Р. к. содержит всю траекторию $\left.f^{t} x\right)$. В этом случае $f^{t} x$ наз. р ек у р р е т т н й т р а е к т о р и е й.



Те о ре м а Би рк г оф а: если пространство $S$ полное (напр., $S=\mathbb{R}^{n}$ ), то 1) для того чтобы точка была рекуррентной, необходимо и достаточно, чтобы замыкание ее траектории было минимальныл множеством; 2) для того чтобы существовала Р. т., достаточно, чтоУстойчивость по Лагранжу).



Р. т. устойчива по Лагранжу и по Пуассону (см. Устойчивость по Пуассону). Почти периодическая (в частности, неподвижная или периодическая) точка динамич. системы рекуррентна. Вообще, всякая точка строго эргодической динамич. системы рекуррентна, но сужение динамич. системы на замыкание рекуррентной траектории (минимальное множество) может не быть строго эргодической динамич. системой (п р и ме р М а р к о в а, см. [2]).



Лит.: [1] Б и р к г о ф Д. Д., Динамические системы, пер. с англ., М., 1941; [2] Н е м ы ц к и й В. В., С т е п а но в В. В., Качественная теория дифференциальных уравнений, 2 изд.,

М.- Л., 1949.



РЕКУ РРЕНТНАЯ ФОРМУЛА - то же, что рекур-



РЕКУ РРЕНТНАЯ ФУНКЦИЯ - Функция, являющаяся рекуррентной точкой сдвигов динамич. системы. Эквивалентное определение: функция $\varphi: \mathbb{R} \rightarrow S$, где $S-$ метрич. пространство, наз. рекуррентной, если она имеет предкомпактное множество значений, равномерно непрерывна и для всякой последовательности чисел $t_{k} \in \mathbb{R}$ такой, что существует предел



\[

\tilde{\varphi}(t)=\lim _{k \rightarrow \infty} \varphi\left(t_{k}+t\right)

\]



(пред е л в ком пактн о о ткры той то пол о г и и, т. е. равномерный на каждом отрезке), найбы существовала точка, устойчивая по Лагранжу (см. рентное соотношение.





\end{document}